% day1_seg4_prob/prob_dist.tex
\begin{frame}
  \onslide<1->{
	\frametitle{Statistics} 
	\framesubtitle{Probability Distributions and Densities}
}
\begin{itemize}
  \onslide<2->{
\item Probability Distributions  
}
  \onslide<3->{
\item Probability Density Functions  
}
\end{itemize}
\end{frame}

\begin{frame}
  \onslide<1->{
	\frametitle{Statistics} 
	\framesubtitle{Probability Distributions}
}
\begin{itemize}
  \onslide<2->{
 \item Every random variable $X$ has an associated probability \textbf{\textit{distribution}} function
 defined by
\[F_X(x)=P\{X^{-1}[-\infty,x)\}=P\set{\omega\in\Omega:X(\omega)< x}\]\\
Also referred to as a \textbf{\textit{Cumulative Distribution Function (c.d.f.)}}
}
  \onslide<3->{
 \item The value of the distribution function at a real number $x$ is the
 probability that the value of $X$ does not exceed $x$.
}
  \onslide<4->{
 \item We can simplify the notation by writing $\set{X< x}$
 instead of $\set{\omega\in\Omega : X(\omega)< x}$.
}
  \onslide<5->{
 \item Thus, $F_X(x)=P\set{X< x}$.
}
\end{itemize}
\end{frame}

\begin{frame}
  \onslide<1->{
	\frametitle{Statistics} 
	\framesubtitle{Probability Distributions}
}
\begin{itemize}
  \onslide<2->{
 \item Distributions are nondecreasing:  $x\leq y\Rightarrow F(x)\leq
 F(y)$.
}
  \onslide<3->{
 \item Distributions are always continuous from the left:  \\
 $F(x)\rightarrow F(a)$ as $x\uparrow a$.
}
  \onslide<4->{
 \item The value of a distribution is always between 0 and 1
 inclusively, and $\lim_{x\downarrow -\infty}F(x)=0$ and $\lim_{x\uparrow \infty}F(x)=1$
}
\end{itemize}
\end{frame}

\begin{frame}
  \onslide<1->{
	\frametitle{Statistics} 
	\framesubtitle{Probability Distributions: Three Flavors}
}
\begin{itemize}
  \onslide<2->{
\item Continuous (no jumps)
 }
  \onslide<3->{
 \item Piecewise continuous (some jumps)
}
  \onslide<4->{
 \item Discrete (only jumps and flat).  $X$ is discrete if and only if $X$
 takes on only finitely many, or at most countably infinite,
 distinct values.  Note that a simple measurable function is a
 discrete random variable.
}
\end{itemize}
\end{frame}

\begin{frame}
  \onslide<1->{
	\frametitle{Statistics} 
	\framesubtitle{Some Discrete Distributions}
}
\begin{itemize}
  \onslide<2->{
 \item Discrete Uniform Distribution:  Let $A=\set{a_1,\ldots,a_n}$ be a set of $n$ distinct real numbers and assume that $a_1<a_2<\ldots<a_n$.  
 Then a random variable $X$ is said to have a discrete uniform
 distribution if \\
 $F(x)= \frac{1}{n}\cdot \abs{A\cap [-\infty, x)}$.
}
  \onslide<3->{
 \item Let $\Omega$ be any sample space and $A\subseteq \Omega$.
 Let $X(\omega)=\chi_A(\omega)$.  Recall that
 $\chi_A(\omega)=
  \begin{cases}
    1 & \text{if } \omega\in A, \\
    0 & \text{otherwise}.
  \end{cases}
$ \\
Then $F(x)=
  \begin{cases}
    0   & \text{if } x<0, \\
    1-P(A) & \text{if } 0\leq x<1, \\
    1   &  \text{if } x\geq 1.
  \end{cases}
$
}
\end{itemize}
\end{frame}

\begin{frame}
  \onslide<1->{
	\frametitle{Statistics} 
	\framesubtitle{Continuous Distributions:  Probability Density Functions}
}
\begin{itemize}
  \onslide<2->{
 \item $X$ is said to have a continuous distribution if there
 exists a function $f:\Real\rightarrow \Real$ such that for each
 $x\in \Real$ \\
\begin{equation}\label{pdf1}
   F(x)=\int_{-\infty}^{x}f(u)du.
\end{equation}
}
  \onslide<3->{
 \item $f$ is called a \textit{\textbf{probability density
 function}}, or p.d.f, for short.
 \item The p.d.f of a random variable, if it exists, is not
 unique.  Why?
}
\end{itemize}
\end{frame}

\begin{frame}
  \onslide<1->{
	\frametitle{Statistics} 
	\framesubtitle{Some Important Distributions}
}
\begin{itemize}
  \onslide<2->{
 \item The Gaussian or normal distribution with p.d.f:
\begin{equation}\label{gauss_distr}
 f(x)=\frac{1}{\sigma\sqrt{2\pi}}\exp\left[-\frac{(x-\mu)^2}{2\sigma^2}\right]
\end{equation}
}
  \onslide<3->{
 \item The uniform distribution on $[a,b]$ with distribution function:
\begin{equation}\label{uniform_distr}
  F(x)=
  \begin{cases}
    0 & \text{if } x< a, \\
    \frac{x-a}{b-a} & \text{if } a\leq x< b,\\
    1 & \text{otherwise}.
  \end{cases}
\end{equation}
}
\end{itemize}
\end{frame}
